\section{Beschreibung des Unternehmens}
\label{sec:unternehmen}

Alamos Gold Inc.\ ist ein kanadischer Goldproduzent mittlerer Gr\"o{\ss}e mit
drei f\"ordernden Betrieben in Nordamerika: dem Island-Gold-Distrikt und der
Young-Davidson-Mine im Norden Ontarios sowie dem Mulatos-Distrikt im
mexikanischen Bundesstaat Sonora.\quelleMDA{2025}{7}\quelleMerken{s4mdaxcacfxh} Die Aktie notiert an der
Toronto Stock Exchange und an der New York Stock Exchange unter dem K\"urzel
AGI.

Dieser Teil ist beschreibend. Wertungen stehen ab
Teil~\ref{sec:kurs}; im laufenden Text sind sie blau gesetzt.

\subsection{Handelnde Personen}

Dem Verwaltungsrat geh\"oren \DirektorenGesamt{} zur Wahl vorgeschlagene
Mitglieder an, von denen \DirektorenUnabh{} als unabh\"angig im Sinne des
\NIOffenlegung{} gelten.\quelleCirc{70}\quelleMerken{s4circxha} Nicht
unabh\"angig ist allein der Vorstandsvorsitzende. Vorsitzender des
Verwaltungsrats ist J.~Robert~S.\ Prichard; Pr\"asident und
Vorstandsvorsitzender ist John~A.\ McCluskey, der das Unternehmen seit seiner
Gr\"undung f\"uhrt.\quelleCirc{15}\quelleMerken{s4circxbf}

Die \"ubrigen Mitglieder sind Alexander Christopher, Elaine Ellingham, David
Fleck, Serafino Tony Giardini, Claire Kennedy, Chana Martineau, Richard
McCreary, Monique Mercier und Shaun Usmar.\quelleCirc{19}\quelleMerken{s4circxbj} Die
Gesch\"aftsleitung besteht neben McCluskey aus Gregory Fisher (Finanzen),
Luc Guimond (Betrieb), Chris Bostwick (Technik) und Scott~R.\,G.\ Parsons
(Exploration).\quelleCirc{54}\quelleMerken{s4circxfe}

\bk{Bemerkenswert ist die fachliche Besetzung des Verwaltungsrats. Vier der
\DirektorenGesamt{} Mitglieder f\"uhren einen geowissenschaftlichen oder
ingenieurwissenschaftlichen Titel, und der Vorsitz des Technik- und
Nachhaltigkeitsausschusses liegt bei einem Wirtschaftspr\"ufer mit
Bergbauerfahrung. Bei einem Unternehmen, dessen Wert an der Umwandlung von
Ressourcen in Reserven h\"angt, ist das kein Zierrat.}

Die f\"unf Aussch\"usse werden gef\"uhrt von Claire Kennedy (Pr\"ufung),
Monique Mercier (Personal), David Fleck (Unternehmensf\"uhrung und
Nominierung), Serafino Tony Giardini (Technik und Nachhaltigkeit) und
J.~Robert~S.\ Prichard (\"Offentliche Angelegenheiten).\quelleCirc{78}\quelleMerken{s4circxhi}

\subsection{Eigent\"umerstruktur}
\label{sec:eigentuemer}

Zum Stichtag des Einberufungsschreibens waren \num{\AktienUmlauf} Aktien
ausgegeben und im Umlauf, jede mit einer Stimme.\quelleCirc{5}\quelleMerken{s4circxf}
Genau ein Aktion\"ar h\"alt mehr als ein Zehntel: Van~Eck Associates
Corporation mit \num{\VanEckAktien} Aktien oder
\Pct{\VanEckAnteil}.\quelleCirc{6}\quelleMerken{s4circxg} Das Unternehmen erkl\"art
ausdr\"ucklich, dass ihm kein weiterer Halter dieser Gr\"o{\ss}enordnung
bekannt ist.

\begin{table}[htbp]
\centering
\caption{Eigent\"umerstruktur zum Stichtag des Einberufungsschreibens.}
\label{tab:eigentuemer}
\begin{tabular}{lrr}
\toprule
Halter & Aktien & Anteil\\
\midrule
Van~Eck Associates Corp. & \num{\VanEckAktien} & \Pct{\VanEckAnteil}\\
\"Ubriger Streubesitz & \num{\StreubesitzAktien} & \Pct{\StreubesitzRest}\\
\midrule
Ausgegeben und im Umlauf & \num{\AktienUmlauf} & \Pct{\AnteilGesamt}\\
\bottomrule
\end{tabular}
\end{table}

\FloatBarrier

\bk{Die Struktur ist die eines Streubesitzunternehmens ohne Ankeraktion\"ar.
Das hat zwei Seiten: Es gibt keinen Gro{\ss}aktion\"ar, der die Strategie
gegen den Kapitalmarkt durchsetzen k\"onnte, und es gibt auch keinen, dessen
Ausstieg den Kurs bestimmt. F\"ur eine Anlage ist die Abwesenheit eines
kontrollierenden Eigent\"umers hier eher ein Vorteil, weil die Aktie ohne
Beherrschungsabschlag gehandelt wird.}

\subsection{Beurteilung durch die Kapitalm\"arkte}

Zum Schlusskurs vom \KursStand{} von \USDje{\Kurs} und den oben genannten
umlaufenden Aktien betr\"agt die B\"orsenbewertung
\MrdUSD{\Marktkapitalisierung}. Sie liegt damit beim
\num[round-mode=places,round-precision=2]{\KursBuchwert}-fachen des
bilanziellen Eigenkapitals je Aktie von \USDje{\BuchwertJeAktie}
(Abschnitt~\ref{sec:bilanz}) und beim
\num[round-mode=places,round-precision=1]{\KursGewinn}-fachen des
Konzernergebnisses des Gesch\"aftsjahres 2025.

Der Kursverlauf der Jahre 2021 bis 2026 und seine Erkl\"arung stehen in
Teil~\ref{sec:kurs}.

\subsection{Analystenabdeckung und Konsens}
\label{sec:konsens}

Kein Unternehmen ver\"offentlicht Ratings oder Kursziele \"uber sich selbst;
solche Angaben stehen in keinem der ausgewerteten Prim\"ardokumente und sind
dort auch nicht zu erwarten. Sie sind jedoch beschaffbar, und ihr Fehlen
w\"are deshalb keine L\"ucke, sondern eine unfertige Suche.

Nach der Erhebung von S\&P~Global verfolgen \num{\KonsensZahl} H\"auser die
Aktie; das aggregierte Votum lautet \emph{Strong Buy}, das mittlere Kursziel
\USDje{\KonsensZiel}, die Spanne reicht von \USDje{\KonsensZielTief} bis
\USDje{\KonsensZielHoch}.\quelleWeb{Stock Analysis, Alamos Gold (AGI) Stock
Forecast}{quelle:konsens}{16.09.2026}

\begin{table}[htbp]
\centering
\caption{Analystenkonsens. \emph{Sekund\"arquelle} -- diese Zeilen tragen
keine gedruckte Seite und w\"urden einer Pr\"ufung gegen ein Prim\"ardokument
nicht standhalten.}
\label{tab:konsens}
\begin{tabular}{lr}
\toprule
Angabe & Wert\\
\midrule
Zahl der sch\"atzenden H\"auser & \num{\KonsensZahl}\\
Aggregiertes Votum & Strong Buy\\
Mittleres Kursziel & \USDje{\KonsensZiel}\\
Niedrigstes Kursziel & \USDje{\KonsensZielTief}\\
H\"ochstes Kursziel & \USDje{\KonsensZielHoch}\\
Ausgewiesener Aufschlag & \Pct{\KonsensAufschlag}\\
\midrule
Zur\"uckgerechneter Bezugskurs & \USDje{\KonsensBezugskurs}\\
Schlusskurs derselben Seite & \USDje{\KonsensKurs}\\
\bottomrule
\end{tabular}
\end{table}

\FloatBarrier

Die beiden letzten Zeilen sind die Aktualit\"atspr\"ufung. Ein Kursziel, das
mit einem Aufschlag ver\"offentlicht wird, verr\"at den Kurs, auf den es
gerechnet wurde: Kursziel geteilt durch eins plus Aufschlag. Das Ergebnis
\USDje{\KonsensBezugskurs} stimmt mit dem auf derselben Seite genannten
Schlusskurs \"uberein. Der Konsens ist damit aktuell und wird nicht als
veraltete Angabe verworfen.

Wie dieses Votum zum Urteil dieses Dokuments steht, behandelt
Abschnitt~\ref{sec:begruendung}.

\subsection{Nicht verf\"ugbare Angaben}
\label{sec:luecken}

Zwei Angaben liegen nicht vor, und beide werden benannt statt gesch\"atzt.

Die Zahl der die Aktie verfolgenden H\"auser wird von den Anbietern
unterschiedlich angegeben; die hier verwendete Erhebung nennt
\num{\KonsensZahl}, andere Zusammenstellungen nennen weniger. Keine der
Quellen erkl\"art die Abweichung, und das Unternehmen selbst nennt keine
Zahl. Der Unterschied ist eine Tatsache \"uber die Daten und wird nicht
gegl\"attet.

Die Aufteilung des Streubesitzes unterhalb der Zehnprozentschwelle
ver\"offentlicht das Unternehmen nicht. Tabelle~\ref{tab:eigentuemer} weist
sie deshalb als einen Posten aus und teilt sie nicht nach gesch\"atzten
Anteilen auf.
